\chapter{Computational Geometry in Three Dimensions}
\label{ch:polyhedral_geometry}

The convex hull of a finite set of points in $\mathbb{R}^3$ is the smallest convex
polyhedron containing them. It is the three-dimensional analogue of the
planar convex hull and is a foundation for collision detection, visibility,
solid modeling, meshing, and geometric optimization.

\section{Polyhedral Representation}
\label{sec:polyhedral_representation}

A polyhedron can be represented by its vertices, edges, and polygonal faces,
or by half-spaces of the form $a_i^{\mathsf T}x\leq b_i$. A valid boundary
representation records face cycles and the adjacency between faces and edges.
For a closed connected convex polyhedron, Euler's formula is
\[
  V-E+F=2.
\]
The boundary orientation must be consistent so that all outward normals point
in the same direction.

\section{Three-Dimensional Predicates}
\label{sec:three_dimensional_predicates}

The signed volume test
\[
\operatorname{orient3d}(A,B,C,D)=
\det\begin{pmatrix}(B-A)^{\mathsf T}\\(C-A)^{\mathsf T}\\(D-A)^{\mathsf T}\end{pmatrix}
\]
determines which side of the oriented plane $ABC$ contains $D$. A positive
value indicates one side, a negative value the other, and zero coplanarity.
Robust evaluation is essential: an incorrect sign can create a non-manifold
polyhedron or reverse a face during hull construction.

\section{Incremental Three-Dimensional Hulls}
\label{sec:incremental_3d_hull}

An incremental hull algorithm starts with a non-degenerate tetrahedron and
inserts points one at a time. For each point, it finds the faces visible from
that point, removes those faces, extracts the horizon cycle between visible
and invisible faces, and connects the new point to the horizon edges.

\begin{algorithm}[htbp]
\caption{Incremental Three-Dimensional Convex Hull}
\label{alg:incremental_3d_hull}
\begin{algorithmic}[1]
\Function{ConvexHull3D}{$P$}
  \State Select four non-coplanar points and initialize tetrahedron $H$
  \For{each remaining point $p$}
    \State $V\gets$ faces of $H$ visible from $p$
    \If{$V$ is non-empty}
      \State $R\gets$ boundary edges between visible and invisible faces
      \State Remove $V$ from $H$
      \For{each edge $(u,v)$ in $R$}
        \State Add face $(u,v,p)$ with consistent outward orientation
      \EndFor
    \EndIf
  \EndFor
  \State \Return $H$
\EndFunction
\end{algorithmic}
\end{algorithm}

The visible-face search can be accelerated with a conflict graph storing the
points visible from each face. With suitable randomized ordering, expected
running time is $O(n\log n+k)$ for $n$ input points and $k$ output features;
the worst case is output-sensitive and can be quadratic for poorly ordered
inputs.

\section{Degeneracies and Validation}
\label{sec:polyhedral_degeneracies}

Implementations must define behavior for coplanar points, duplicate points,
collinear subsets, cospherical configurations, and points lying on existing
faces. Symbolic perturbation or exact predicates can make the combinatorial
result deterministic. A resulting hull should be checked for positive face
orientation, edge incidence, connectedness, Euler's formula, and containment
of every input point.

\section{Applications}

Three-dimensional convex hulls provide collision proxies, bounding volumes,
support mappings for GJK, visibility boundaries, polyhedral approximations,
and input domains for tetrahedral and hexahedral mesh generation.

